% =====================================================================
% LAMPIRAN D: SKEMA BASIS DATA SUPABASE LENGKAP
% =====================================================================
\chapter{Skema Basis Data Supabase Lengkap}
\label{app:db-schema-detail}

Lampiran ini menampilkan rincian kolom seluruh tabel basis data Supabase, daftar \textit{enum} pendukung, dan kebijakan \textit{Row Level Security} (RLS) yang diringkas pada Bagian~\ref{sec:db-schema}. Pemisahan dilakukan agar pembahasan rancangan sistem pada Bab~\ref{ch:metode} tetap padat sementara skema penuh tetap dapat diaudit secara mendetail.

\section*{D.1 Tabel \texttt{profiles}}

Tabel \texttt{profiles} memuat informasi profil pengguna yang terhubung satu-ke-satu dengan tabel sistem \texttt{auth.users} milik Supabase Auth. Atribut utama dirangkum pada Tabel~\ref{tab:profiles}.

\begin{table}[htbp]
\centering
\caption{Atribut tabel \texttt{profiles}.}
\label{tab:profiles}
\small
\begin{tabular}{p{3.4cm} p{2.6cm} p{6.5cm}}
\toprule
\textbf{Kolom} & \textbf{Tipe} & \textbf{Deskripsi} \\
\midrule
id & uuid (PK) & Mengacu ke \texttt{auth.users.id}. \\
full\_name & text & Nama lengkap pengguna. \\
email & text & Alamat email pengguna. \\
phone\_number & text & Nomor telepon kontak. \\
profile\_photo\_url & text & URL foto profil pada Supabase Storage. \\
is\_moderator & bool & Penanda peran moderator. \\
role & user\_role & \textit{Enum} peran (warga, operator, moderator). \\
assigned\_agency & issue\_category & \textit{Enum} instansi penugasan operator. \\
created\_at & timestamptz & Waktu pembuatan profil. \\
updated\_at & timestamptz & Waktu perubahan terakhir. \\
\bottomrule
\end{tabular}
\end{table}

\section*{D.2 Tabel \texttt{reports}}

Tabel \texttt{reports} adalah tabel utama yang menyimpan setiap laporan warga beserta hasil klasifikasi AI dan metadata lokasi. Atribut utama dirangkum pada Tabel~\ref{tab:reports}.

\begin{table}[htbp]
\centering
\caption{Atribut tabel \texttt{reports}.}
\label{tab:reports}
\small
\begin{tabular}{p{3.4cm} p{2.6cm} p{6.5cm}}
\toprule
\textbf{Kolom} & \textbf{Tipe} & \textbf{Deskripsi} \\
\midrule
id & uuid (PK) & Identitas unik laporan. \\
user\_id & uuid (FK) & Mengacu ke \texttt{profiles.id}. \\
reporter\_name & text & Nama pelapor saat mengirim. \\
reporter\_email & text & Email pelapor saat mengirim. \\
category & issue\_category & \textit{Enum} kategori dinas akhir. \\
ai\_prediction & text & Kategori dinas yang diprediksi AI. \\
ai\_confidence & numeric & Tingkat keyakinan model (0..1). \\
ai\_probabilities & jsonb & Skor probabilitas untuk setiap dari sembilan instansi target. \\
ai\_model\_* & text & Metadata model server-side, mencakup nama dan versi model. \\
assigned\_* & text/float8 & \textit{Proyeksi-baca} penugasan instansi terkini (sumber: \texttt{report\_assignments}). \\
image\_url & text & \textit{Proyeksi-baca} referensi foto awal (sumber: \texttt{report\_evidence}). \\
description & text & Narasi laporan dalam Bahasa Indonesia. \\
latitude & float8 & Koordinat lintang lokasi kejadian. \\
longitude & float8 & Koordinat bujur lokasi kejadian. \\
address & text & Alamat hasil \textit{reverse geocoding}. \\
visibility & report\_visibility & \textit{Enum} visibilitas (\textit{public}/\textit{private}). \\
duplicate\_of & uuid (FK) & Menunjuk laporan asal bila laporan ini duplikat. \\
status & report\_status & \textit{Enum} status (\textit{submitted}, \textit{verified}, \textit{in\_progress}, \textit{resolved}, \textit{rejected}). \\
verified\_by & uuid (FK) & Operator yang memverifikasi laporan. \\
verified\_at & timestamptz & Waktu verifikasi. \\
rejection\_reason & text & Alasan penolakan bila status \textit{rejected}. \\
resolution\_* & text/\allowbreak{}uuid/\allowbreak{}timestamptz & \textit{Proyeksi-baca} bukti/metadata penyelesaian (sumber: \texttt{report\_evidence}). \\
created\_at & timestamptz & Waktu laporan dikirim. \\
updated\_at & timestamptz & Waktu perubahan terakhir. \\
\bottomrule
\end{tabular}
\end{table}

Kolom \texttt{assigned\_*}, \texttt{image\_url}, dan \texttt{resolution\_*} pada tabel \texttt{reports} merupakan \textit{proyeksi-baca}: nilainya disinkronkan otomatis oleh \textit{trigger} \texttt{SECURITY DEFINER} dari tabel anak \texttt{report\_assignments} dan \texttt{report\_evidence} yang menjadi sumber kebenaran. Proyeksi ini dipertahankan agar klien Flutter yang membaca melalui Supabase \textit{Realtime} (\textit{stream} satu tabel) tetap menerima satu baris datar.

\section*{D.3 Tabel \texttt{report\_evidence}}

Tabel \texttt{report\_evidence} menyimpan banyak berkas bukti per laporan (foto awal, foto penyelesaian, atau bukti tambahan). Referensi penyimpanan dicatat sebagai pasangan \textit{bucket} dan \texttt{object\_path}, bukan URL final, sehingga URL akses dapat dibuat sebagai \textit{signed URL} berumur singkat saat dibutuhkan. Atribut utama dirangkum pada Tabel~\ref{tab:report-evidence}.

\begin{table}[htbp]
\centering
\caption{Atribut tabel \texttt{report\_evidence}.}
\label{tab:report-evidence}
\small
\begin{tabular}{p{3.4cm} p{2.6cm} p{6.5cm}}
\toprule
\textbf{Kolom} & \textbf{Tipe} & \textbf{Deskripsi} \\
\midrule
id & uuid (PK) & Identitas unik berkas bukti. \\
report\_id & uuid (FK) & Mengacu ke \texttt{reports.id}. \\
kind & evidence\_kind & Jenis bukti (\textit{initial}, \textit{resolution}, \textit{additional}). \\
bucket & text & Nama \textit{bucket} Supabase Storage. \\
object\_path & text & Lintasan objek di dalam \textit{bucket}. \\
content\_type & text & Tipe MIME berkas. \\
uploaded\_by & uuid (FK) & Pengguna yang mengunggah. \\
created\_at & timestamptz & Waktu unggah. \\
\bottomrule
\end{tabular}
\end{table}

\section*{D.4 Tabel \texttt{report\_assignments}}

Tabel \texttt{report\_assignments} merekam riwayat penugasan instansi untuk tiap laporan, sehingga penugasan-ulang (\textit{reassignment}) dan eskalasi dapat ditelusuri. Hanya satu baris berstatus \textit{active} yang berlaku per laporan (dijamin indeks unik parsial). Atribut utama dirangkum pada Tabel~\ref{tab:report-assignments}.

\begin{table}[htbp]
\centering
\caption{Atribut tabel \texttt{report\_assignments}.}
\label{tab:report-assignments}
\small
\begin{tabular}{p{3.4cm} p{2.6cm} p{6.5cm}}
\toprule
\textbf{Kolom} & \textbf{Tipe} & \textbf{Deskripsi} \\
\midrule
id & uuid (PK) & Identitas unik penugasan. \\
report\_id & uuid (FK) & Mengacu ke \texttt{reports.id}. \\
agency\_id & text & Referensi logis ke \texttt{agencies.id}. \\
agency\_name & text & Nama instansi tujuan. \\
agency\_category & issue\_category & Kategori instansi tujuan. \\
distance\_meters & float8 & Jarak instansi dari koordinat laporan. \\
routing\_method & text & Metode penentuan instansi terdekat. \\
status & assignment\_status & \textit{Enum} status (\textit{active}, \textit{reassigned}, \textit{escalated}, \textit{superseded}). \\
reason & text & Alasan penugasan/eskalasi. \\
assigned\_by & uuid (FK) & Aktor penugasan (kosong bila otomatis). \\
created\_at & timestamptz & Waktu penugasan. \\
\bottomrule
\end{tabular}
\end{table}

\section*{D.5 Tabel \texttt{agencies}}

Tabel \texttt{agencies} adalah registri instansi tujuan yang dirujuk oleh \texttt{report\_assignments}, menggantikan medan teks instansi yang sebelumnya tersimpan datar pada \texttt{reports}. Atribut utama dirangkum pada Tabel~\ref{tab:agencies}.

\begin{table}[htbp]
\centering
\caption{Atribut tabel \texttt{agencies}.}
\label{tab:agencies}
\small
\begin{tabular}{p{3.4cm} p{2.6cm} p{6.5cm}}
\toprule
\textbf{Kolom} & \textbf{Tipe} & \textbf{Deskripsi} \\
\midrule
id & text (PK) & Identitas unit instansi. \\
name & text & Nama instansi. \\
category & issue\_category & Kategori dinas instansi. \\
latitude & float8 & Koordinat lintang unit (opsional). \\
longitude & float8 & Koordinat bujur unit (opsional). \\
is\_active & bool & Penanda instansi aktif. \\
created\_at / updated\_at & timestamptz & Waktu pembuatan dan perubahan. \\
\bottomrule
\end{tabular}
\end{table}

\section*{D.6 Tabel \texttt{report\_history}}

Tabel \texttt{report\_history} mencatat setiap perubahan status pada laporan agar jejak audit (\textit{audit trail}) dapat ditelusuri dari awal pengiriman hingga penyelesaian. Atribut utama dirangkum pada Tabel~\ref{tab:report-history}.

\begin{table}[htbp]
\centering
\caption{Atribut tabel \texttt{report\_history}.}
\label{tab:report-history}
\small
\begin{tabular}{p{3.4cm} p{2.6cm} p{6.5cm}}
\toprule
\textbf{Kolom} & \textbf{Tipe} & \textbf{Deskripsi} \\
\midrule
id & uuid (PK) & Identitas unik catatan riwayat. \\
report\_id & uuid (FK) & Mengacu ke \texttt{reports.id}. \\
status & report\_status & Status baru pada langkah ini. \\
note & text & Catatan operator atau sistem. \\
updated\_by & uuid & Pengguna yang memicu perubahan. \\
created\_at & timestamptz & Waktu perubahan. \\
\bottomrule
\end{tabular}
\end{table}

\section*{D.7 Tabel \texttt{report\_sla}}

Tabel \texttt{report\_sla} menyimpan tenggat layanan per laporan untuk memantau pemenuhan SLA; kebijakan SLA tingkat instansi diserahkan sebagai pengembangan lanjutan. Atribut utama dirangkum pada Tabel~\ref{tab:report-sla}.

\begin{table}[htbp]
\centering
\caption{Atribut tabel \texttt{report\_sla}.}
\label{tab:report-sla}
\small
\begin{tabular}{p{3.8cm} p{2.4cm} p{6.3cm}}
\toprule
\textbf{Kolom} & \textbf{Tipe} & \textbf{Deskripsi} \\
\midrule
report\_id & uuid (PK, FK) & Mengacu ke \texttt{reports.id}. \\
first\_response\_due\_at & timestamptz & Tenggat respons pertama. \\
resolution\_due\_at & timestamptz & Tenggat penyelesaian. \\
first\_response\_at & timestamptz & Waktu respons pertama tercapai. \\
breached\_at & timestamptz & Waktu SLA terlanggar. \\
created\_at / updated\_at & timestamptz & Waktu pembuatan dan perubahan. \\
\bottomrule
\end{tabular}
\end{table}

\section*{D.8 \textit{Enum} Pendukung}

Skema mendefinisikan enam \textit{enum} kustom PostgreSQL yang dipakai bersama oleh tabel-tabel di atas, yaitu:

\begin{itemize}
    \item \texttt{issue\_category}: kategori dinas dengan nilai sesuai sembilan kategori target pada Tabel~\ref{tab:kategori-target}.
    \item \texttt{report\_status}: status siklus hidup laporan, yaitu \textit{submitted}, \textit{verified}, \textit{in\_progress}, \textit{resolved}, \textit{rejected}.
    \item \texttt{user\_role}: peran pengguna, yaitu \textit{citizen}, \textit{agency\_admin}, \textit{super\_admin}.
    \item \texttt{evidence\_kind}: jenis bukti, yaitu \textit{initial}, \textit{resolution}, \textit{additional}.
    \item \texttt{assignment\_status}: status penugasan, yaitu \textit{active}, \textit{reassigned}, \textit{escalated}, \textit{superseded}.
    \item \texttt{report\_visibility}: visibilitas laporan, yaitu \textit{public}, \textit{private}.
\end{itemize}

\section*{D.9 \textit{Row Level Security} dan Penyimpanan}

Kebijakan \textit{Row Level Security} (RLS) diaktifkan pada seluruh tabel agar akses data terkontrol pada level baris. Beberapa kebijakan utama yang dirancang adalah:

\begin{itemize}
    \item Setiap pelapor hanya dapat membaca dan memperbarui baris \texttt{reports} miliknya sendiri (\(\texttt{auth.uid()} = \texttt{user\_id}\)).
    \item Petugas dengan \texttt{role = 'agency\_admin'} dapat membaca dan memperbarui laporan pada \texttt{assigned\_agency} yang sesuai dengan \texttt{category}.
    \item Admin dengan \texttt{role = 'super\_admin'} memiliki akses lintas instansi, termasuk koreksi penugasan; penulisan tabel \texttt{agencies} dibatasi pada peran ini.
    \item Tabel anak \texttt{report\_evidence} dan \texttt{report\_assignments} hanya dapat dibaca oleh pengguna terautentikasi; penulisannya dilakukan eksklusif melalui \textit{trigger} \texttt{SECURITY DEFINER} sehingga tidak ada kebijakan \texttt{INSERT} langsung.
    \item Tabel \texttt{report\_history} hanya ditulis melalui RPC/pemicu basis data saat status \texttt{reports} berubah.
\end{itemize}

Berkas disimpan pada \textit{bucket} privat \texttt{report-images} dan \texttt{profile-photos}. Kebijakan unggah membatasi penulisan objek ke prefiks \texttt{\{user\_id\}/} milik pengguna, sedangkan pembacaan dilakukan melalui \textit{signed URL} berumur singkat yang dibuat atas dasar kebijakan \texttt{SELECT} untuk pengguna terautentikasi; basis data menyimpan \textit{bucket} dan \texttt{object\_path}, bukan URL final. Skema \texttt{storage} bawaan Supabase ditunjukkan pada Gambar~\ref{fig:storage-schema}.

\begin{figure}[htbp]
    \centering
    \includegraphics[width=0.95\textwidth]{supabase/storage_schema.png}
    \caption{Skema \texttt{storage} yang dikelola Supabase. Aplikasi hanya merujuk secara logis tabel \texttt{objects} (melalui \texttt{bucket\_id} serta \texttt{name}/\texttt{object\_path}) dan \texttt{buckets}; tabel lain seperti \texttt{s3\_multipart\_uploads} dan \texttt{vector\_indexes} merupakan internal Supabase Storage dan tidak digunakan aplikasi.}
    \label{fig:storage-schema}
\end{figure}
